\subsection{扩张定理}

设 \(\mathfrak{g} = \mathfrak{h} + \mathfrak{g}'\) 是李代数，它是理想 \(\mathfrak{g}'\) 与子代数 \(\mathfrak{h}\) 的直和；令 U 为 \(\mathfrak{g}\) 的包络代数，令 \(\mathrm{U}' \subset \mathrm{U}\) 为包含于 U 的、\(\mathfrak{g}'\) 的包络代数。存在唯一的 \(\mathfrak{g}\)-模结构，定义在 \(\mathrm{U}'\) 上，使得：(α) 对于 \(x \in \mathfrak{g}'\) 和 \(u \in \mathrm{U}'\)，有 \(x_{\mathrm{U}'}u = -ux\)；(β) 对于 \(x \in \mathfrak{h}\) 和 \(u \in \mathrm{U}'\)，有 \(x_{\mathrm{U}'}u = xu - ux\)（后一个元素确实属于 \(\mathrm{U}'\)，因为由 \(x\) 定义的 U 的内导子保持 \(\mathfrak{g}'\)，因而也保持 \(\mathrm{U}'\) 稳定）。由条件（α）和（β）唯一确定映射 \(x \mapsto x_{\mathrm{U}'}\)，它把 \(\mathfrak{g}\) 映入 \(\mathcal{L}_{\mathrm{K}}(\mathrm{U}')\)。所以只需验证 \([x, y]_{\mathrm{U}'} = [x_{\mathrm{U}'}, y_{\mathrm{U}'}]\)；只需考虑以下情形：

\begin{enumerate}
\def\labelenumi{(\arabic{enumi})}
\tightlist
\item
  若 \(x \in \mathfrak{g}', y \in \mathfrak{g}'\)，则
\end{enumerate}

\[
[ x, y ] _ {\mathrm{U} ^ {\prime}} u = - u (x y - y x) = \left(x _ {\mathrm{U} ^ {\prime}} y _ {\mathrm{U} ^ {\prime}} - y _ {\mathrm{U} ^ {\prime}} x _ {\mathrm{U} ^ {\prime}}\right) u;
\]

\begin{enumerate}
\def\labelenumi{(\arabic{enumi})}
\setcounter{enumi}{1}
\tightlist
\item
  若 \(x \in \mathfrak{h}, y \in \mathfrak{g}'\)，则
\end{enumerate}

\[
\begin{array}{r l} {[ x, y ] _ {\mathrm{U}'} u = - u (x y - y x) = x (- u y) - (- u y) x + (x u - u x) y} \\ & {= (x _ {\mathrm{U}'} y _ {\mathrm{U}'} - y _ {\mathrm{U}'} x _ {\mathrm{U}'}) u;} \end{array}
\]

\begin{enumerate}
\def\labelenumi{(\arabic{enumi})}
\setcounter{enumi}{2}
\tightlist
\item
  若 \(x \in \mathfrak{h}, y \in \mathfrak{h}\)，则 \([x, y]_{\mathrm{U}'}\) 和 \([x_{\mathrm{U}'}, y_{\mathrm{U}'}]\) 是 \(\mathrm{U}'\) 的两个导子；它们在 \(\mathfrak{g}'\) 上的限制分别与 \(\operatorname{ad}_{\mathfrak{g}}[x, y]\) 和 \([\operatorname{ad}_{\mathfrak{g}}x, \operatorname{ad}_{\mathfrak{g}}y]\) 的限制相同，因而这两个导子相等。
\end{enumerate}

还要考虑对偶表示 \(x \mapsto -^{t}x_{\mathrm{U}'}\)，它是 \(\mathfrak{g}\) 在 \(\mathrm{U}'^*\) 上的表示。对于 \(x \in \mathfrak{g}', -^{t}x_{\mathrm{U}'}\)，它是右乘 \(x\) 的 \(\mathrm{U}'\) 的转置；相应的 \(\mathrm{U}'\) 的表示是 \(\mathrm{U}'\) 的余正则表示。

定义 1. 设 g 是李代数，g' 是其子代数，\(\rho'\) 是 g' 在 V' 上的表示。若 \(\rho\) 是 g 在 V 上的表示，且存在从 g'-模 V' 到 g'-模 V 的单射同态，则称它为 \(\rho'\) 向 g 的扩张。我们也说 g-模 V 是 g'-模 V' 的扩张。

若 \(\rho'\) 是有限维的，且 \(\mathfrak{g}'\) 是 \(\mathfrak{g}\) 的可解理想，则要存在有限维扩张，必要条件是 \([\mathfrak{g},\mathfrak{g}']\) 包含于 \(\rho'\) 的最大幂零理想中（\(\S 5\)，第 3 号，定理 1）。

定理 1（Zassenhaus）。设 \(\mathfrak{g} = \mathfrak{g}' + \mathfrak{h}\) 是李代数，它是理想 \(\mathfrak{g}'\) 与子代数 \(\mathfrak{h}\) 的直和，且 \(\rho'\) 是 \(\mathfrak{g}'\) 的有限维表示，其最大幂零理想包含 \([\mathfrak{h},\mathfrak{g}']\)。

\begin{enumerate}
\def\labelenumi{(\alph{enumi})}
\item
  存在 \(\rho\)，它是 \(\rho'\) 向 \(\mathfrak{g}\) 的有限维扩张，且其最大幂零理想包含 \(\rho'\) 的最大幂零理想。
\item
  若对所有 \(x \in \mathfrak{h}\)，\(\mathfrak{g}'\) 上的 \(\operatorname{ad}_{\mathfrak{g}}x\) 的限制是幂零的，则可选择 \(\rho\)，使得 \(\rho\) 的最大幂零理想还包含 \(\mathfrak{h}\)。
\end{enumerate}

令 \(U'\) 为 \(g'\) 的包络代数。假设 \(U'\) 和 \(U'^*\) 具有本节开头定义的 \(g\)-模结构。

令 \(I \subset U^{\prime}\) 为 \(\rho^{\prime}\)（视为 \(U^{\prime}\) 的表示）的核。它是 \(U^{\prime}\) 的有限余维双边理想。\(C(\rho^{\prime})\)（参见第 1 号）是 \(U^{\prime*}\) 中与 I 正交的子空间。令 S 为 \(U^{\prime*}\) 的由 \(C(\rho^{\prime})\) 生成的 g-子模。

\[
\mathrm{U}^{\prime *} \supset \mathrm{S} \supset \mathrm{C}(\rho')
\]

现在证明 S 在 K 上有限维。令 \(V'\) 为 \(\rho'\) 所作用的空间，且 \(V' = V_{0}' \supset V_{1}' \supset \cdots \supset V_{d}' = \{0\}\) 是 \(g'\)-模 \(V'\) 的 Jordan--Hölder 序列。令 \(\rho_{i}'\) 为 \(g'\) 在 \(V_{i-1}'/V_{i}'\) 上由 \(\rho' (1 \leqslant i \leqslant d)\) 诱导的表示。令 \(I' \subset U'\) 为各个 \(\rho_{i}'\) 的核的交（把这些表示视为 \(U'\) 的表示）。

\[
\mathrm{I} ^ {\prime d} \subset \mathrm{I} \subset \mathrm{I} ^ {\prime}
\]

并且 \(\mathrm{I}'\cap \mathfrak{g}'\) 是 \(\rho'\) 的最大幂零理想。根据 § 2，第 6 号，命题 6 的推论，\(\mathrm{I}^{\prime d}\) 在 \(\mathrm{U}'\) 中具有有限余维。若 \(x\in \mathfrak{h}\)，则 \(u\mapsto xu - ux\) 是 \(\mathrm{U}'\) 的导子；它把 \(\mathfrak{g}'\) 映入 \([\mathfrak{h},\mathfrak{g}']\subset \mathrm{I}'\)，因而把 \(\mathrm{U}'\) 映入 \(\mathrm{I}'\)，继而把 \(\mathrm{I}^{\prime d}\) 映入 \(\mathrm{I}^{\prime d}\)。另一方面，显然 \(\mathrm{I}^{\prime d}\) 是 \(\mathfrak{g}'\)-模 \(\mathrm{U}'\) 的子模。因此 \(\mathrm{I}^{\prime d}\) 是 \(\mathfrak{g}\)-模 \(\mathrm{U}'\) 的子模。\(\mathrm{I}^{\prime d}\) 在 \(\mathrm{U}^*\) 中的正交补是有限维的 \(\mathfrak{g}\)-子模，包含 \(\mathrm{C}(\rho')\)，因而也包含 S。这说明 S 在 K 上有限维。对于 \(x\in I'\cap \mathfrak{g}', x^d\)，它显然包含于 \(\mathfrak{g}\)-模 \(\mathrm{U}' / \mathrm{I}^{\prime d}\) 的湮没子中，因而也包含于 \(\mathfrak{g}\)-模 S 的湮没子中。

第 1 号说明，\(\mathfrak{g}'\)-模 \(\mathbf{V}'\) 同构于 \(\mathfrak{g}'\)-模 \((\mathrm{C}(\rho'))^n\) 的一个子模。因此，\(\mathfrak{g}\)-模 \(\mathrm{S}^n\) 给出 \(\rho\)，它是 \(\rho'\) 向 \(\mathfrak{g}\) 的有限维扩张。此外，\(\rho(x)\) 对于 \(x \in \mathrm{I}' \cap \mathfrak{g}'\) 是幂零的；由于 \(\mathrm{I}' \cap \mathfrak{g}'\) 是 \(\mathfrak{g}\) 的理想（因为按假设它包含 \([\mathfrak{h},\mathfrak{g}']\)），可知 \(\mathrm{I}'\cap \mathfrak{g}'\) 包含于 \(\rho\) 的最大幂零理想中。因此（a）得证。

最后，假设对所有 \(x \in \mathfrak{h}\)，\(\mathfrak{g}'\) 上的 \(\mathrm{ad}_{\mathfrak{g}}x\) 的限制是幂零的。由于 \(\mathfrak{h}\) 中的元素通过导子作用于 \(\mathrm{U}'\)，所以对所有 \(u \in \mathrm{U}'\) 和所有 \(x \in \mathfrak{h}\)，存在整数 \(e\) 使 \((x_{\mathfrak{U}'})^e.u = 0\)；因此由 \(x_{\mathfrak{U}'}\) 在 \(\mathrm{U}' / \mathrm{I}^{\prime d}\) 和 S 上诱导的自同态是幂零的。于是 \(\rho(x)\) 对所有 \(x \in \mathfrak{h}\) 都是幂零的。我们先前已知，对于 \(x \in \mathrm{I}' \cap \mathfrak{g}'\) 也成立。由于 \(\mathrm{I}' \cap \mathfrak{g}'\) 是 \(\mathfrak{g}'\) 的理想且包含 \([\mathfrak{h}, \mathfrak{g}']\)，所以 \(\mathfrak{h} + (\mathrm{I}' \cap \mathfrak{g}')\) 也是 \(\mathfrak{g}\) 的理想。定理 1 的（b）由下面的引理得到：

引理 1. 设 \(\mathfrak{g} = \mathfrak{g}' + \mathfrak{h}\) 是李代数，是理想 \(\mathfrak{g}'\) 与子代数 \(\mathfrak{h}\) 的和。\(\sigma\) 是 \(\mathfrak{g}\) 的有限维表示。假设 \(\sigma(x)\) 对所有 \(x \in \mathfrak{g}'\) 和 \(x \in \mathfrak{h}\) 都是幂零的。则 \(\sigma(x)\) 对所有 \(x \in \mathfrak{g}\) 都是幂零的。

将 g 换成 \(\sigma\) 的核的商后，可设 \(\sigma\) 是忠实的。于是 \(\mathfrak{g}'\) 和 \(\mathfrak{h}\) 是幂零的，因此 \(\mathfrak{g}\)（它是 \(\mathfrak{h}\) 的某个商被 \(\mathfrak{g}'\) 扩张）是可解的。于是 \(\mathfrak{h}\) 和 \(\mathfrak{g}'\) 都包含于 \(\sigma\) 的最大幂零理想中（§ 5，第 3 号，定理 1 的推论 6）。

关于定理 1 的改进，参见练习 4。

